%=============================================================================
% PATENT 15 — RESEARCH SUMMARY
% Comprehensive overview tying together all deliverables
%=============================================================================
\documentclass[11pt,a4paper]{article}

\usepackage[margin=1in]{geometry}
\usepackage{amsmath,amssymb}
\usepackage{siunitx}
\usepackage{booktabs}
\usepackage{longtable}
\usepackage{hyperref}
\usepackage{xcolor}
\usepackage{enumitem}
\usepackage{tcolorbox}

\title{\textbf{Patent 15 Research Summary}\\[6pt]
\large Methods, Systems, and Apparatus for Artificial Generation,
Control, and Application of Scalar Density Fields\\[6pt]
\normalsize US Provisional Patent 63/868,073 (Aug 21, 2025)}

\author{DFD Research Programme --- Patent 15 Agent}
\date{March 2026}

\begin{document}
\maketitle

%=============================================================================
\section{Mission Summary}
%=============================================================================

Patent 15 is the \textbf{master methods patent} for the entire DFD
technology portfolio.  It covers the fundamental apparatus for
generating and controlling the scalar density field $\psi$---the
enabling technology for all other DFD patents (propulsion,
beamforming, energy transfer, navigation, sensing, computing, etc.).

This research output comprises four documents:

\begin{enumerate}[nosep]
\item \textbf{LaTeX Paper}: ``Artificial Generation and Control of
Scalar Density Fields'' --- the foundational technical paper
establishing how to build a $\psi$ generator, with complete
signal-level predictions and noise budgets.

\item \textbf{Engineering Design}: Complete deployment-ready
specification for the Laboratory Scalar Density Field Generator
(LSDFG), including chamber specs, source configurations, sensor
layouts, control architecture, safety systems, cost estimates,
and schedules.

\item \textbf{New Patent Proposals}: Four new patent opportunities
identified during the engineering design process.

\item \textbf{This Summary}: Overview and key findings.
\end{enumerate}

%=============================================================================
\section{Key Technical Findings}
%=============================================================================

\subsection{The $\lambda$ Parameter Is the Gate}

Everything depends on the EM--$\psi$ back-reaction parameter
$\lambda$.  The core DFD postulates specify how $\psi$ affects EM
propagation ($n = e^\psi$), but $\lambda$ addresses the inverse:
can EM fields modify $\psi$?

\begin{itemize}[nosep]
\item $\lambda = 1$: EM cannot generate $\psi$.  Laboratory
$\psi$ generation is impossible through EM channels.
\item $|\lambda - 1| \neq 0$: EM can pump $\psi$ modes.
The proposed apparatus can generate, detect, and control $\psi$.
\end{itemize}

Current status: existing cavity stability provides an accidental
bound $|\lambda - 1| \lesssim 3 \times 10^{-5}$.  The proposed
apparatus improves this by up to \textbf{nine orders of magnitude}
to $|\lambda - 1| \sim 10^{-14}$.

\subsection{Geometry Transparency Must Be Broken}

Symmetric cavities in pure eigenmodes cannot drive $\psi$ modes
because the EM stress tensor integral vanishes ($G = 0$).  The
critical insight is that \textbf{TE+TM superposition} with
controlled relative phase breaks this transparency, restoring
$G \neq 0$ with cross-mode overlap parameter $\eta_\times \sim 0.3$.

This is the single most important technical finding: without
geometry breaking, no amount of stored energy produces a $\psi$
field.  With it, even modest stored energies ($\sim 10\;\si{kJ}$)
can produce detectable signals.

\subsection{The Design Law Governs Everything}

The master design equation:
\begin{equation}
|\lambda - 1|_{\rm min} = \frac{\pi\gamma_\psi}{c_s\,U_0\,m}
\frac{A_\psi^2}{\kappa_{\rm eff}\,A_{\rm cav,tot}}
\end{equation}
shows that sensitivity improves with:
\begin{itemize}[nosep]
\item Higher stored energy $U_0$ (linear)
\item Deeper modulation $m$ (linear)
\item Larger cavity aperture $A_{\rm cav,tot}$ (linear)
\item Smaller tube area $A_\psi$ (quadratic)
\item Lower $\psi$-mode damping $\gamma_\psi$ (linear)
\end{itemize}

The quadratic dependence on tube area is notable: reducing
$A_\psi$ by a factor of 3 improves sensitivity by a factor of 9.

\subsection{Multi-Modal Detection Provides Cross-Validation}

A genuine $\psi$ perturbation produces correlated signals in
three independent domains:
\begin{itemize}[nosep]
\item Optical: refractive index change $\delta n = n\,\delta\psi$
\item Atomic: frequency shift $\delta\nu/\nu = K_A\,\delta\psi$
\item Inertial: acceleration $a = (c^2/2)\,\nabla(\delta\psi)$
\end{itemize}

Any signal appearing in only one domain is a systematic error.
This three-way cross-check provides extremely high confidence
in any claimed detection.

%=============================================================================
\section{Sensitivity Projections}
%=============================================================================

\begin{longtable}{p{0.18\textwidth}p{0.18\textwidth}p{0.22\textwidth}p{0.18\textwidth}p{0.14\textwidth}}
\toprule
\textbf{Phase} & \textbf{Timeline} & \textbf{$\Delta\psi_{\rm min}$}
& \textbf{$|\lambda{-}1|_{\rm min}$} & \textbf{Cost} \\
\midrule
\endhead
Phase 1 & Year 1 & $\sim 10^{-15}$ & $\sim 10^{-9}$ & \$5.3M \\
Phase 2 & Year 2 & $\sim 10^{-17}$ & $\sim 10^{-11}$ & +\$1.5M \\
Phase 3 & Year 3+ & $\sim 10^{-19}$ & $\sim 10^{-14}$ & +\$3M \\
\bottomrule
\end{longtable}

\textbf{Improvement over current bounds}: Phase 1 alone improves on
existing accidental constraints by 4 orders of magnitude.

%=============================================================================
\section{Prior Art Context}
%=============================================================================

\subsection{Key Related Experimental Work}

\begin{itemize}[nosep]
\item \textbf{SRF cavities}: TESLA-type cavities at DESY, Fermilab,
and JLab routinely achieve $Q_0 > 10^{10}$.  The proposed apparatus
leverages this mature technology base.

\item \textbf{Analog gravity}: Acoustic black holes in BECs
(Steinhauer 2016), optical analogs (Philbin 2008), and dynamical
Casimir effect (Wilson 2011) demonstrate that effective spacetime
metrics can be created in laboratory systems.  DFD goes further:
it claims the vacuum \emph{is} the medium.

\item \textbf{Precision clocks}: The BACON network, PTB Yb$^+$
E3/E2, and JILA Sr clocks provide the measurement infrastructure.
These experiments already constrain DFD parameters and validate
the clock-comparison technique.

\item \textbf{Interferometry}: LIGO, LISA Pathfinder, and
the Fermilab Holometer demonstrate the interferometric sensitivity
needed.  The proposed MZI requires ``only''
$10^{-10}\;\si{rad/\sqrt{Hz}}$, far less demanding than LIGO.

\item \textbf{Superconducting gravimeters}: Commercial instruments
achieve $10^{-12}\;\si{m/s^2/\sqrt{Hz}}$, sufficient for Phase 1
operation.
\end{itemize}

\subsection{Why This Has Not Been Done Before}

Three factors explain why the $\psi$-generation experiment has not
been attempted:

\begin{enumerate}[nosep]
\item \textbf{Geometry transparency}: Pure-mode symmetric cavities
have $G = 0$.  The driven channel is invisible in standard
operation.  The TE+TM breaking technique requires deliberate
design.

\item \textbf{$2\omega$ suppression}: Standard cavity metrology
actively suppresses $2\omega$ components as unwanted AM sidebands.
The very signal we seek is being filtered out in existing
experiments.

\item \textbf{No theoretical motivation}: Without DFD's prediction
of the $\lambda$ parameter and the specific mode equation,
there is no reason to look for parametric $\psi$ generation
in EM cavities.
\end{enumerate}

%=============================================================================
\section{New Patent Opportunities}
%=============================================================================

Four new patents identified:

\begin{tcolorbox}[colback=blue!3, colframe=blue!50!black]
\begin{description}[nosep]
\item[A: Dual-Mode Cavity $\psi$-Pump] --- The specific TE+TM
technique breaking geometry transparency.
\textbf{Priority: Highest.}

\item[B: Multi-Modal Sensor Fusion] --- Kalman filter estimation
of $\psi(\mathbf{x},t)$ from optical + atomic + inertial data.
\textbf{Priority: Medium.}

\item[C: $\psi$-Mode Spectroscopy] --- Swept-frequency discovery
of scalar field resonances with $2\omega$ lock-in.
\textbf{Priority: Medium.}

\item[D: Safety Interlock System] --- Hardware-enforced $\psi$-field
limits with microsecond response.
\textbf{Priority: High.}
\end{description}
\end{tcolorbox}

Total estimated filing costs: \$90K (all four through PCT).

%=============================================================================
\section{Recommendations}
%=============================================================================

\begin{enumerate}
\item \textbf{File Proposal A immediately}: The dual-mode cavity
pump is the core enabling invention.  Provisional filing should
proceed within 3 months.

\item \textbf{Begin Phase 1 procurement}: Long-lead items (SRF
cavities, He refrigerator) have 6--9 month lead times.

\item \textbf{Engage SRF facility}: Partner with DESY, Fermilab,
or JLab for cavity fabrication and initial testing.

\item \textbf{Develop safety standards}: No regulatory framework
exists for artificial gravitational fields.  Proposal D establishes
the foundation; engage safety consultants early.

\item \textbf{Publish the LaTeX paper}: Establishing priority in
the open literature strengthens the patent position and invites
collaboration.

\item \textbf{Target existing anomalies}: Review SRF cavity
operational logs for unexplained $2\omega$ signals that may
have been dismissed as technical noise.  These could be the
first evidence of $|\lambda - 1| \neq 0$.
\end{enumerate}

%=============================================================================
\section{Deliverables Checklist}
%=============================================================================

\begin{longtable}{p{0.05\textwidth}p{0.55\textwidth}p{0.15\textwidth}p{0.15\textwidth}}
\toprule
& \textbf{Deliverable} & \textbf{File} & \textbf{Status} \\
\midrule
\endhead
1 & Full LaTeX paper on $\psi$-field generation &
\texttt{Paper.tex} & Complete \\
2 & Engineering design specification &
\texttt{Design.tex} & Complete \\
3 & New patent proposals (4 proposals) &
\texttt{Proposals.tex} & Complete \\
4 & Research summary (this document) &
\texttt{Summary.tex} & Complete \\
\bottomrule
\end{longtable}

All deliverables written to:\\
\texttt{DFD\_Research\_Output/Patent\_15\_Provisional\_Methods\_Systems/}

\end{document}
